% GENERATED FILE. Edit canonical lessons, then run npm run generate:books.
% canonical-source-hash: fnv1a64:3a9050532696176a
% canonical-lessons: MW-C198-daphtar, MW-C198-aphsar, MW-C198-arji, MW-C198-binti, MW-C198-samjhavno

\chapter{At the Office}
\label{ch:mw-at-the-office}
\hlchaptermodality{198}

\begin{chapteropening}
\textbf{By the end of this chapter:} \emph{I can ask for something at an office.}

Complete the last lesson of chapter 198: I can ask for something at an office.
\end{chapteropening}

\section[\texorpdfstring{\mw{दफ्तर} (daphtar)}{दफ्तर (daphtar)}]{\mw{दफ्तर} (daphtar) — an office}

\label{lesson:MW-C198-daphtar}



\begin{quote}
Before the new one: say the Marwadi for a message, then the Marwadi for to rub out.
\end{quote}

\subsection*{You'll want to know: \mw{दफ्तर}}
\textbf{\mw{दफ्तर}} — \emph{daphtar} — \textquotedblleft{}an office\textquotedblright{}.

\begin{quote}
\mw{दफ्तर} \mw{आज} \mw{बंद} \mw{है।}
\end{quote}

\begin{quote}
\emph{daphtar āj band hai.}
\end{quote}

\begin{quote}
The office is closed today.
\end{quote}

A masculine noun.

\begin{cousinweb}
\textbf{\mw{दफ्तर}} (\emph{daphtar}) is usually traced, through Persian, to Arabic \emph{daftar}, \textquotedblleft{}a register\textquotedblright{}.
\end{cousinweb}

\begin{grammarlens}[title={What you've built}]
One more word for reading notices and messages.
\end{grammarlens}

\subsection*{Your turn}
\begin{itemize}
\raggedright
  \item \emph{Say it:} \emph{daphtar}
  \item \emph{Say it:} \emph{daphtar}, once more
  \item \emph{Say it:} \emph{daphtar āj band hai.} — The office is closed today.
  \item \emph{Recall:} say the Marwadi for a message, then the Marwadi for to rub out, then say \emph{daphtar} again
\end{itemize}

\begin{tcolorbox}[breakable,colback=teal!4,colframe=teal!35!black,title={Before you move on}]
What does \mw{दफ्तर} mean? (\textquotedblleft{}An office\textquotedblright{}.) Say it once more.
\end{tcolorbox}
\section[\texorpdfstring{\mw{अफसर} (aphsar)}{अफसर (aphsar)}]{\mw{अफसर} (aphsar) — an officer}

\label{lesson:MW-C198-aphsar}



\begin{quote}
Before the new one: say the Marwadi for to rub out, then the Marwadi for an office.
\end{quote}

\subsection*{You'll want to know: \mw{अफसर}}
\textbf{\mw{अफसर}} — \emph{aphsar} — \textquotedblleft{}an officer\textquotedblright{}.

\begin{quote}
\mw{अफसर} \mw{दफ्तर} \mw{मांय} \mw{कोनी।}
\end{quote}

\begin{quote}
\emph{aphsar daphtar mā̃y konī.}
\end{quote}

\begin{quote}
The officer isn't in the office.
\end{quote}

A masculine or feminine noun, by the person meant.

\begin{cousinweb}
\textbf{\mw{अफसर}} (\emph{aphsar}) is borrowed from English \emph{officer}.
\end{cousinweb}

\begin{grammarlens}[title={What you've built}]
One more word for reading notices and messages.
\end{grammarlens}

\subsection*{Your turn}
\begin{itemize}
\raggedright
  \item \emph{Say it:} \emph{aphsar}
  \item \emph{Say it:} \emph{aphsar}, once more
  \item \emph{Say it:} \emph{aphsar daphtar mā̃y konī.} — The officer isn't in the office.
  \item \emph{Recall:} say the Marwadi for to rub out, then the Marwadi for an office, then say \emph{aphsar} again
\end{itemize}

\begin{tcolorbox}[breakable,colback=teal!4,colframe=teal!35!black,title={Before you move on}]
What does \mw{अफसर} mean? (\textquotedblleft{}An officer\textquotedblright{}.) Say it once more.
\end{tcolorbox}
\section[\texorpdfstring{\mw{अरजी} (arjī)}{अरजी (arjī)}]{\mw{अरजी} (arjī) — an application}

\label{lesson:MW-C198-arji}



\begin{quote}
Before the new one: say the Marwadi for an office, then the Marwadi for an officer.
\end{quote}

\subsection*{You'll want to know: \mw{अरजी}}
\textbf{\mw{अरजी}} — \emph{arjī} — \textquotedblleft{}an application\textquotedblright{}.

\begin{quote}
\mw{हूं} \mw{अरजी} \mw{लिखूं} \mw{हूं।}
\end{quote}

\begin{quote}
\emph{hū̃ arjī likhū̃ hū̃.}
\end{quote}

\begin{quote}
I am writing an application.
\end{quote}

A feminine noun.

\begin{cousinweb}
\textbf{\mw{अरजी}} (\emph{arjī}) is usually traced, through Persian, to Arabic \emph{ʿarḍ}, \textquotedblleft{}a presenting\textquotedblright{}: an application lays a matter before an officer.
\end{cousinweb}

\begin{grammarlens}[title={What you've built}]
One more word for reading notices and messages.
\end{grammarlens}

\subsection*{Your turn}
\begin{itemize}
\raggedright
  \item \emph{Say it:} \emph{arjī}
  \item \emph{Say it:} \emph{arjī}, once more
  \item \emph{Say it:} \emph{hū̃ arjī likhū̃ hū̃.} — I am writing an application.
  \item \emph{Recall:} say the Marwadi for an office, then the Marwadi for an officer, then say \emph{arjī} again
\end{itemize}

\begin{tcolorbox}[breakable,colback=teal!4,colframe=teal!35!black,title={Before you move on}]
What does \mw{अरजी} mean? (\textquotedblleft{}An application\textquotedblright{}.) Say it once more.
\end{tcolorbox}
\section[\texorpdfstring{\mw{बिनती} (bintī)}{बिनती (bintī)}]{\mw{बिनती} (bintī) — a request}

\label{lesson:MW-C198-binti}



\begin{quote}
Before the new one: say the Marwadi for an officer, then the Marwadi for an application.
\end{quote}

\subsection*{You'll want to know: \mw{बिनती}}
\textbf{\mw{बिनती}} — \emph{bintī} — \textquotedblleft{}a request\textquotedblright{}.

\begin{quote}
\mw{म्हारी} \mw{बिनती} \mw{सुणो।}
\end{quote}

\begin{quote}
\emph{mhārī bintī suṇo.}
\end{quote}

\begin{quote}
Please hear my request.
\end{quote}

A feminine noun.

\begin{cousinweb}
\textbf{\mw{बिनती}} (\emph{bintī}) is inherited from Sanskrit \emph{vijñapti}, \textquotedblleft{}a request to a superior\textquotedblright{}; CDIAL lists Old Marwari \emph{bīnatī}, \textquotedblleft{}request\textquotedblright{} (CDIAL 11706).
\end{cousinweb}

\begin{grammarlens}[title={What you've built}]
One more word for reading notices and messages.
\end{grammarlens}

\subsection*{Your turn}
\begin{itemize}
\raggedright
  \item \emph{Say it:} \emph{bintī}
  \item \emph{Say it:} \emph{bintī}, once more
  \item \emph{Say it:} \emph{mhārī bintī suṇo.} — Please hear my request.
  \item \emph{Recall:} say the Marwadi for an officer, then the Marwadi for an application, then say \emph{bintī} again
\end{itemize}

\begin{tcolorbox}[breakable,colback=teal!4,colframe=teal!35!black,title={Before you move on}]
What does \mw{बिनती} mean? (\textquotedblleft{}A request\textquotedblright{}.) Say it once more.
\end{tcolorbox}
\section[\texorpdfstring{\mw{समझावणो} (samjhāvṇo)}{समझावणो (samjhāvṇo)}]{\mw{समझावणो} (samjhāvṇo) — to explain}

\label{lesson:MW-C198-samjhavno}



\begin{quote}
Before the new one: say the Marwadi for an application, then the Marwadi for a request.
\end{quote}

\subsection*{You'll want to know: \mw{समझावणो}}
\textbf{\mw{समझावणो}} — \emph{samjhāvṇo} — \textquotedblleft{}to explain\textquotedblright{}.

\begin{quote}
\mw{अफसर} \mw{नियम} \mw{समझावै} \mw{है।}
\end{quote}

\begin{quote}
\emph{aphsar niyam samjhāvai hai.}
\end{quote}

\begin{quote}
The officer explains the rule.
\end{quote}

Past (I …): \mw{म्हैं} \mw{बात} \mw{समझाई}, \emph{mhaĩ bāt samjhāī}.

\begin{cousinweb}
\textbf{\mw{समझावणो}} (\emph{samjhāvṇo}) is the causative of \textbf{\mw{समझणो}}, \textquotedblleft{}to understand\textquotedblright{}: to make someone understand. Both are inherited from Sanskrit \emph{sambudhyate}, \textquotedblleft{}wakes up, understands\textquotedblright{}; CDIAL lists Old Marwari \emph{samajhāi}, \textquotedblleft{}persuades\textquotedblright{} (CDIAL 12959).
\end{cousinweb}

\begin{grammarlens}[title={What you've built}]
One more word for reading notices and messages.
\end{grammarlens}

\subsection*{Your turn}
\begin{itemize}
\raggedright
  \item \emph{Say it:} \emph{samjhāvṇo}
  \item \emph{Say it:} \emph{samjhāvṇo}, once more
  \item \emph{Say it:} \emph{aphsar niyam samjhāvai hai.} — The officer explains the rule.
  \item \emph{Recall:} say the Marwadi for an application, then the Marwadi for a request, then say \emph{samjhāvṇo} again
\end{itemize}

\begin{tcolorbox}[breakable,colback=teal!4,colframe=teal!35!black,title={Before you move on}]
What does \mw{समझावणो} mean? (\textquotedblleft{}To explain\textquotedblright{}.) Say it once more.
\end{tcolorbox}
